% =====================================================================
\chapter{Naive Bayes Classifier}
% =====================================================================

\labinfo{05\_Naive\_Bayes}{naive\_bayes.py}{UCI SMS Spam Collection: 5,572 text messages (747 spam)}
{Binary text classification (spam or ham)}{Test accuracy 99.0\%, spam F1 0.963, ROC AUC 0.991 ($\alpha=0.25$)}

\section{Dataset}
The \textbf{SMS Spam Collection} (Almeida and Gómez Hidalgo, 2011; UCI) contains 5,572 real English text
messages labelled \emph{ham} (normal, 4,827) or \emph{spam} (747, i.e.\ 13.4\%). A typical spam message
looks like ``WINNER!! As a valued network customer you have been selected to receive a £900 prize
reward! To claim call 0906\dots''. Spam messages are longer on average (about 24 words against 14).

Preparation was deliberately simple: lower-case everything and keep runs of letters, digits and the
\pounds{} and \$ signs as words. No stemming and no stop-word removal. I used a stratified 80/20 split
(4,460 training and 1,114 test messages). The training messages contain 7,735 different words.

Because the classes are unbalanced, accuracy alone would be misleading (always answering ``ham'' already
scores 86.6\%), so I focus on precision, recall and F1 for the spam class.

\section{Code and experiments}
\begin{itemize}
\item \code{tokenize} is a one-line regular expression.
\item \code{MultinomialNaiveBayes.fit} builds the vocabulary, counts words per class and stores
      $\log P(c)$ and the smoothed $\log P(w\mid c)$ for every word. With $\alpha=0$, unseen words get
      $\log 0=-\infty$ on purpose, to show the zero-frequency problem.
\item \code{log\_scores} adds up the log-probabilities of the known words; \code{predict} takes the
      larger score; \code{spam\_score} returns the score difference for the ROC curve.
\item \code{binary\_metrics}, \code{roc\_curve} and \code{auc} compute the evaluation numbers.
\end{itemize}
\noindent Experiments: (1)~5-fold cross-validation over nine values of $\alpha$ from 0 to 5; (2)~test-set
comparison of no smoothing, Laplace ($\alpha=1$) and the CV choice; (3)~the most ``spammy'' words, i.e.\
the largest $\log P(w\mid\text{spam})-\log P(w\mid\text{ham})$.

\section{Results}

\begin{table}[H]
\centering\small
\caption{5-fold cross-validation F1 score (spam class) for different smoothing values.}
\label{tab:nb-alpha}
\begin{tabular}{@{}l c c c c c c c c c@{}}
\toprule
$\alpha$ & 0 & 0.01 & 0.05 & 0.1 & 0.25 & 0.5 & 1 & 2 & 5\\
\midrule
CV F1 & 0.784 & 0.940 & 0.945 & 0.946 & \textbf{0.949} & 0.943 & 0.943 & 0.928 & 0.867\\
\bottomrule
\end{tabular}
\end{table}

\begin{table}[H]
\centering\small
\caption{Test-set results (1,114 messages, spam is the positive class).}
\label{tab:nb-test}
\begin{tabular}{@{}l S[table-format=1.4] S[table-format=1.4] S[table-format=1.4] S[table-format=1.4] S[table-format=1.4] S[table-format=1.4]@{}}
\toprule
\textbf{Model} & {\textbf{Accuracy}} & {\textbf{Precision}} & {\textbf{Recall}} & {\textbf{F1}} & {\textbf{Specificity}} & {\textbf{AUC}}\\
\midrule
No smoothing ($\alpha=0$)  & 0.9551 & 0.9714 & 0.6846 & 0.8031 & 0.9969 & 0.9730\\
Laplace ($\alpha=1$)       & 0.9910 & 0.9860 & 0.9463 & 0.9658 & 0.9979 & 0.9875\\
CV choice ($\alpha=0.25$)  & 0.9901 & 0.9662 & 0.9597 & 0.9630 & 0.9948 & 0.9913\\
\bottomrule
\end{tabular}
\end{table}

\textbf{Smoothing is essential.} Without it, 1,027 of the 1,114 test messages contain at least one word
that has zero probability in one of the classes. Those messages are decided by which class happened to
see the rarest word rather than by their content, and recall falls to 68\%: a third of the spam gets
through. Any $\alpha$ between 0.01 and 1 fixes this.

The CV choice $\alpha=0.25$ and Laplace smoothing are practically tied on the test set; the difference
is two or three messages. $\alpha=0.25$ catches two more spam messages at the cost of three more false
alarms. Both reach about 99\% accuracy, which is far above the 86.6\% of the always-ham baseline.

The learned word statistics are easy to interpret: \emph{claim}, \emph{prize}, \emph{guaranteed},
\emph{awarded}, prices such as \emph{150p} and \emph{\pounds1000}, and \emph{ringtone}/\emph{tones}
are typical of SMS spam. Such an easy explanation is one of the nice things about Naive Bayes.

\begin{figure}[H]
\centering
\includegraphics[width=0.48\linewidth]{\figpath{05\_Naive\_Bayes}{roc_and_alpha.pdf}}\hfill
\includegraphics[width=0.48\linewidth]{\figpath{05\_Naive\_Bayes}{top_spam_words.pdf}}
\caption{ROC and smoothing comparison, with the most spam-indicative words.}
\end{figure}

\section{Source code}
\lstinputlisting[caption={\texttt{05\_Naive\_Bayes/naive\_bayes.py}}]{05_Naive_Bayes/naive_bayes.py}
\clearpage
